\section{Identifiability and code-frequency consistency}
\label{sec:identifiability}

We port \cref{thm:bg-id} to electronic phenotyping. Three results
follow: a glass-ceiling theorem (\cref{thm:glass-ceiling}), an
identifiability theorem with chart review
(\cref{thm:identifiability-chart}), and a code-frequency consistency
corollary (\cref{thm:code-total}).

\subsection{The phenotyping glass ceiling}

A patient with no diagnosis code for condition $Z$ has candidate set
$\{0,1\}$: the codes do not resolve the true state. From code data
alone, the only observable quantity for a single binary code is the
code frequency
\begin{equation}
  q := \Prob(C = 1)
     = \pi \cdot \mathrm{sens} + (1 - \pi) \cdot (1 - \mathrm{spec}),
  \label{eq:code-freq}
\end{equation}
a single scalar constraint on the three unknowns $(\pi, \mathrm{sens},
\mathrm{spec})$.

\begin{theorem}[Glass ceiling for phenotyping]
\label{thm:glass-ceiling}
Let $C$ be a single binary diagnosis code generated under the DGP of
\cref{sec:translation} with unknown prevalence $\pi$, code sensitivity
$\mathrm{sens}$, and code specificity $\mathrm{spec}$. If the coding
mechanism is left unrestricted, then $(\pi, \mathrm{sens},
\mathrm{spec})$ is non-identifiable from code data alone: for any
observed code frequency $q \in (0,1)$ there exists a continuum of
triples $(\pi, \mathrm{sens}, \mathrm{spec})$ producing the same $q$.
\end{theorem}

\begin{proof}
Equation \eqref{eq:code-freq} is one scalar equation in three unknowns
$(\pi, \mathrm{sens}, \mathrm{spec}) \in (0,1)^3$. Fix any target $q
\in (0,1)$. We exhibit a two-dimensional surface of admissible triples
producing the same $q$. Choose $(\pi, \mathrm{spec})$ with $\pi \in
(0,1)$ and $\mathrm{spec}$ in the admissibility region
\begin{equation}
  \max\!\Bigl(0,\ 1 - \tfrac{q}{1-\pi}\Bigr)
  \;\le\; \mathrm{spec} \;\le\;
  \min\!\Bigl(1,\ \tfrac{1-q}{1-\pi}\Bigr),
  \label{eq:admissible}
\end{equation}
and solve \eqref{eq:code-freq} for
\begin{equation}
  \mathrm{sens}(\pi, \mathrm{spec})
    = \frac{q - (1-\pi)(1-\mathrm{spec})}{\pi}.
  \label{eq:glass-sens}
\end{equation}
The two bounds in \eqref{eq:admissible} are exactly the constraints
$\mathrm{sens} \le 1$ (upper bound on $\mathrm{spec}$) and
$\mathrm{sens} \ge 0$ (lower bound): from $\mathrm{sens} \le 1$,
$(1-\pi)(1-\mathrm{spec}) \ge q - \pi$, i.e.
$\mathrm{spec} \le (1-q)/(1-\pi)$. (The earlier statement omitted the
upper bound; without it a triple such as $q=0.5, \pi=0.1,
\mathrm{spec}=0.99$ would give the inadmissible $\mathrm{sens}=4.91$.)
By construction the triple
$(\pi, \mathrm{sens}(\pi, \mathrm{spec}), \mathrm{spec})$
reproduces $q$ exactly. The admissibility region is a two-dimensional
open subset of the unit square whenever it is nonempty (which holds
for every $q \in (0,1)$, e.g. around $\mathrm{spec} = 1-q$ at small
$\pi$), so the solution set is a two-dimensional surface in the unit
cube. Every point of this
surface produces the identical observed code frequency $q$; distinct
parameter triples are therefore observationally indistinguishable from
the code data alone.
\end{proof}

\Cref{thm:glass-ceiling} formalizes a fact that EHR researchers know
informally: the fraction of patients carrying a code is not the
prevalence. Setting $\hat\pi_{\text{code}} = q$ implicitly assumes
$\mathrm{sens} = \mathrm{spec} = 1$, the one corner of the solution
surface where the code is a perfect proxy. Every other admissible
coding model gives a different prevalence consistent with the same
data. Code-only prevalence is biased not because the estimator is
crude but because the parameter is non-identifiable.

This non-identifiability is the EHR-coding instance of a general
phenomenon in partial identification: when a binary outcome is
measured with unknown error, its marginal distribution is only
set-identified. \citet{horowitz1995identification} bound population
parameters under contaminated and corrupted sampling, and
\citet{molinari2008partial} gives sharp identification regions for the
distribution of a misclassified variable. \Cref{thm:glass-ceiling}
specializes that picture to a single diagnosis code with unknown
sensitivity and specificity: rather than bound a functional, we exhibit
the full two-dimensional solution surface in
$(\pi,\mathrm{sens},\mathrm{spec})$ space and show every point on it is
observationally equivalent, which is what makes the chart-review
singletons of \cref{thm:identifiability-chart} the identifying device
that the partial-identification literature would call an external
restriction.

\subsection{Identifiability with chart review}

The role of chart review is to add singleton-candidate-set rows to
$\tilde{C}$, restoring rank. A chart-reviewed patient has an
adjudicated true state $\Y_i$; the pair $(\Y_i, C_i)$ is a fully
observed draw from the coding mechanism.

\begin{theorem}[Identifiability with chart review]
\label{thm:identifiability-chart}
Suppose a chart-reviewed subsample of $m$ patients is adjudicated, and
the subsample contains at least one patient with $\Y = 1$ and at least
one with $\Y = 0$. If
\begin{enumerate}[(i)]
\item the coding mechanism belongs to a parametric family identifiable
  from the joint distribution of $(\Y, S, C)$ on the subsample, and
\item the conditional severity distribution $\Prob(S \mid \Y)$ on the
  subsample equals $\Prob(S \mid \Y)$ on the cohort
  (\emph{case-mix representativeness}: this holds trivially when the
  subsample is drawn iid from the cohort, and violations are
  quantified in \cref{cor:casemix}),
\end{enumerate}
then the coding-model parameters are identified from the subsample, and
the cohort prevalence $\pi$ is identified from the full-cohort code
frequency $q$ via \eqref{eq:code-freq}.
\end{theorem}

\begin{proof}[Sketch]
The argument is two-stage, paralleling the spike-in identifiability
theorem of \citet{towell2026scrnacoarsening}. Stage one: on the
chart-reviewed subsample both $\Y_i$ and $C_i$ are observed, so the
conditional code probabilities $\Prob(C = 1 \mid \Y = 1, S)$ and
$\Prob(C = 1 \mid \Y = 0)$ are estimated by a standard regression of
$C$ on $\Y$ and $S$; condition (i) supplies enough distinct design
points to identify the coding parameters, and hence $\mathrm{sens}$ and
$\mathrm{spec}$. Stage two: with $\mathrm{sens}$ and $\mathrm{spec}$
identified, \eqref{eq:code-freq} becomes one equation in the single
unknown $\pi$, solved as
\begin{equation}
  \hat\pi
    = \frac{q - (1 - \widehat{\mathrm{spec}})}
            {\widehat{\mathrm{sens}} - (1 - \widehat{\mathrm{spec}})},
  \label{eq:deconvolve}
\end{equation}
which is well-defined whenever $\widehat{\mathrm{sens}} + \widehat{\mathrm{spec}}
> 1$, the condition that the code is informative about $\Y$. Equation
\eqref{eq:deconvolve} is the \citet{rogan1978estimating} prevalence
correction estimator with $\mathrm{sens}$ and $\mathrm{spec}$ supplied
by the chart-reviewed subsample instead of being assumed known. The
contribution of this paper at this step is structural (placing the
estimator inside the masked-cause framework and exhibiting the
chart-reviewed singleton as the identifying mechanism), not a new
estimator. Identifiability follows by standard regularity. Condition (ii) is the
representativeness requirement: if it fails, stage two uses subsample
sensitivity where cohort sensitivity is needed, and $\hat\pi$ inherits
that error. This residual bias is bounded in \cref{sec:methodology}.
\end{proof}

\Cref{thm:identifiability-chart} is the phenotyping analog of the
spike-in theorem of \citet{towell2026scrnacoarsening}. A chart-reviewed
patient plays exactly the role that an ERCC spike-in transcript plays
there: a singleton candidate set whose truth is known by construction,
supplying the identifying rows that code data alone cannot. The two
qualifiers deserve emphasis. Condition (i) requires the chart-reviewed
subsample to cover both $\Y = 0$ and $\Y = 1$ patients; reviewing only
code-positive patients identifies sensitivity but not specificity.
Condition (ii) is the case-mix representativeness assumption, the
direct analog of the ERCC-endogenous capture-efficiency match; it is
the principal practical limit and is examined in \cref{sec:methodology}.

\subsection{Code-frequency consistency}

A fitted phenotype model has a structural property: it reproduces the
marginal observed code frequency exactly, regardless of the prevalence
it reports.

\begin{theorem}[Code-frequency consistency]
\label{thm:code-total}
Let $(\hat\pi, \widehat{\mathrm{sens}}, \widehat{\mathrm{spec}})$ be any
interior maximum-likelihood estimate of the phenotype model for a
single binary code $C$. Then the fitted code frequency matches the
empirical code frequency exactly,
\begin{equation}
  \hat\pi \cdot \widehat{\mathrm{sens}}
    + (1 - \hat\pi)\cdot(1 - \widehat{\mathrm{spec}})
    = \bar{C},
  \label{eq:code-total}
\end{equation}
where $\bar{C}$ is the empirical fraction of code-positive patients.
\end{theorem}

\begin{proof}
The marginal log-likelihood of a single binary code depends on the
parameters only through the implied code frequency $q(\pi,
\mathrm{sens}, \mathrm{spec})$ of \eqref{eq:code-freq}. The data
contribute a Bernoulli log-likelihood proportional to $\bar{C}\log q +
(1 - \bar{C})\log(1 - q)$, whose unique stationary point in $q \in
(0,1)$ is $q = \bar{C}$. Any interior MLE in the regime $\mathrm{sens} + \mathrm{spec} > 1$,
i.e.\ where the code is informative about $\Y$ (so the gradient of $q$
with respect to $(\pi, \mathrm{sens}, \mathrm{spec})$ is nonzero on the
parameter interior), therefore satisfies $q(\hat\pi,
\widehat{\mathrm{sens}}, \widehat{\mathrm{spec}}) = \bar{C}$, which is
\eqref{eq:code-total}. This is the binary, code-frequency form of the
cell-total consistency identity of \citet[\S 3]{towell2026scrnacoarsening}:
at an interior MLE the fitted marginal matches the empirical marginal
exactly. The multi-code extension follows by the same argument applied
to the joint code-frequency vector; we omit the algebra here.
\end{proof}

\Cref{thm:code-total} has a sharp practical consequence. Every point on
the glass-ceiling solution surface of \cref{thm:glass-ceiling}
satisfies \eqref{eq:code-total}: a model reporting $\hat\pi = 0.05$ and
a model reporting $\hat\pi = 0.20$ can both reproduce the observed code
frequency perfectly. A marginal-fit check, comparing predicted to
observed code frequency, therefore cannot detect prevalence bias. The
diagnostic that practitioners reach for, ``does the model recover the
observed coding rate,'' is uninformative about the quantity of
interest. Diagnosing prevalence bias requires the singleton candidate
sets of \cref{thm:identifiability-chart}: chart review, anchor
observations, or an external gold-standard registry. The same limit was
identified for scRNA-seq in \citet{towell2026scrnacoarsening} and for
spatial deconvolution in \citet{towell2026spatialcoarsening}.
Code-frequency consistency (\cref{thm:code-total}) is a special case of
the general coarsened-data consistency theorem of
\citet{towell2026synthesis}, in which the singleton-candidate-set
device used here, chart review, appears as one instance of a
domain-independent rank condition and singleton-restoration mechanism.
